\section{Verification-Aware Solver Selection}
\label{sec:formulation}

\subsection{Equation Workload}

Let $q\in\mathcal{Q}$ denote one solve request. A request contains an equation
family, dimensions and sparsity, coefficients or callback handles, initial state,
tolerances, available compute backends, and an original state $s_q$. The equation
family defines an authoritative residual or defect operator
$\mathcal{R}_q(y)$ and a scaled acceptance threshold $\tau_q$. For dynamic
systems, $y$ also contains time, derivatives, mode, and event metadata.

An expert $e\in\mathcal{E}$ may be a complete classical solver, a preconditioner,
a warm start, a learned operator, a device kernel, a corrector, or a verifier. Its
output is therefore typed by a role $\rho(e)$ rather than assumed to be a final
solution. A plan
\begin{equation}
  P=(e_{1:k},c,g,t)
\end{equation}
contains an ordered candidate cascade $e_{1:k}$, optional corrector $c$, mandatory
gate $g$, and terminal numerical solver $t$. A gate is a numerical acceptance test
derived from the original equation. If a stage is rejected, execution continues from
the original request state to the next numerical path.

\subsection{Complete Verified Cost}

Let $R_j$ denote the event that cascade stage $j$ is reached and let $A_j$
denote acceptance at that stage. Thus $R_1$ always occurs and
$R_{j+1}=R_j\cap\neg A_j$. For stage $j$, let $K_j$ include candidate,
transfer, correction, and gate cost; let $C_0$ include assessment, routing, and
shared setup; and let $C_t$ include the terminal solver and its final gate. The
conditional complete-cost model is
\begin{align}
\mathbb{E}[C(P,q)\mid\phi(q)] ={}& C_0(q)
 + \sum_{j=1}^{k}\Pr(R_j\mid\phi(q))\bar K_j(q) \nonumber\\
&+\Pr(R_{k+1}\mid\phi(q))\bar C_t(q),
\label{eq:complete-cost}
\end{align}
where $\bar K_j=\mathbb{E}[K_j\mid R_j,\phi(q)]$ and
$\bar C_t=\mathbb{E}[C_t\mid R_{k+1},\phi(q)]$.
This reach-weighted form makes rejection and numerical continuation explicit. A microsecond
candidate is not beneficial if transfer, correction, verification, or its
probability of reaching the terminal solver dominates the saved classical work.

The ideal router seeks
\begin{equation}
 P^*(q)=\arg\min_{P\in\mathcal{P}(q)}
 \mathbb{E}[C(P,q)\mid \phi(q)]
\label{eq:router-objective}
\end{equation}
subject to capability, memory, artifact, and role constraints. Here $\phi(q)$ is a
feature vector derived without running all candidate experts. A hindsight per-
scenario oracle is used only as an analysis reference; its search cost and query
winner labels are excluded from deployable routing. In the evaluated profiles, stage
costs and reach frequencies come from disjoint training/calibration scenarios, either
as action-level constants or as functions of $\phi(q)$; held-out outcomes are not
available to the online decision.

The objective also yields an inspectable ordering rule. Let
$p_j=\Pr(A_j\mid R_j,\phi(q))>0$ and assume that two eligible stages have
order-invariant conditional costs and acceptance probabilities. Placing stage $a$
before stage $b$ costs $\bar K_a+(1-p_a)\bar K_b$ before their common suffix; the
reverse order costs $\bar K_b+(1-p_b)\bar K_a$. Therefore, $a$ should precede $b$
if and only if
\begin{equation}
 \frac{\bar K_a}{p_a}\leq\frac{\bar K_b}{p_b}.
\label{eq:cascade-index}
\end{equation}
Repeated adjacent exchanges show that sorting by $\bar K_j/p_j$ minimizes the
expected cost of a fixed eligible cascade under these assumptions, consistent with
prior independent-test ordering results~\cite{berend2014tests}. This result
does not choose the globally optimal subset when stages change one another's state
or when a contract fixes terminal-stage order.

Correction effort is itself a routing decision. Let $\mathcal{B}_e$ be the
calibrated budgets available for expert $e$, and let action $a=(e,b)$ have
conditional attempt cost $K_a$ and acceptance probability $p_a$. At most one
budget may be selected for each expert. After sorting actions by $K_a/p_a$, the
optimal suffix value for action index $i$, used-expert set $S$, and remaining
capacity $r$ is
\begin{equation}
\begin{aligned}
 V(i,S,r)=\min\{&V(i+1,S,r),\\
 &K_i+(1-p_i)V(i+1,S\cup\{e_i\},r-1)\},
\end{aligned}
\label{eq:joint-route-budget}
\end{equation}
with the take branch disabled when $e_i\in S$, and terminal value $C_t$ when
$r=0$ or no action remains. The exact dynamic program jointly chooses expert, budget,
and order under a bounded action model. The implementation accepts either static
calibrated $(K_a,p_a)$ or request-conditioned predictions
$(\widehat K_a(\phi),\widehat p_a(\phi))$; the original-equation gate remains outside
the optimizer.

For adjacent transition effects, let $M_{b,a}>0$ multiply $K_a$ after action $b$,
with $M_{\bot,a}=1$. For used experts $S$, define
\begin{equation}
\begin{aligned}
 W(S,b)=\min\{C_t,\min_{a:e_a\notin S}\{&K_aM_{b,a}\,+\\
 &(1-p_a)W(S\cup\{e_a\},a)\}\},
\end{aligned}
\label{eq:interaction-route-budget}
\end{equation}
with the inner minimum disabled when $|S|=k$.

\begin{proposition}[Finite cascade exactness and hardness]
\label{prop:exact-finite-cascade}
\label{prop:exact-interaction-cascade}
\label{prop:interaction-hardness}
With finite nonnegative costs, positive history-independent acceptance, capacity $k$,
and one action per expert: (i) when $M_{b,a}=1$, sorting by $K_a/p_a$ and
Eq.~\eqref{eq:joint-route-budget} are globally exact in $O(n2^m)$ worst-case scalar
states/transitions; (ii) with immediately previous-action multipliers,
$W(\varnothing,\bot)$ is globally exact with $O(n^2 2^m)$ transitions; and
(iii) the rational interaction-aware decision problem is NP-complete even for
$K_a=1$, $p_a=1/2$, $k=n$, $C_t=5$, and $M_{b,a}\in\{1,2\}$.
\end{proposition}

\begin{proof}
Adjacent exchange plus backward induction proves (i); Bellman induction on remaining
experts proves (ii). For (iii), reduce directed Hamiltonian path with one action per
vertex and multipliers one on edges and two otherwise. Appending an action costs at
most $4.5<C_t$, so every optimum uses all vertices, and cost is at most
$\sum_{j=0}^{n-1}2^{-j}+5\,2^{-n}$ exactly for a Hamiltonian path. Membership in NP
follows by evaluating a cascade; complete details are in the supplementary artifact.
\end{proof}

The 64-bit implementations support at most 63 experts and reject on state-cap
overflow; exactness requires the cap to admit every reached state.

\subsection{Numerical Acceptance and Continuation}

The numerical contract is
\begin{equation}
\begin{aligned}
 \operatorname{return}(y)&\Longrightarrow g_q(y)=\mathrm{accept},\\
 g_q(y)&\equiv
 \|D_q^{-1}\mathcal{R}_q(y)\|_\infty\le \tau_q.
\end{aligned}
\label{eq:gate}
\end{equation}
with family-specific additions for backward error, integration defect,
constraints, complementarity, event guards, and DAE hidden residuals. $D_q$
contains absolute, relative, and nominal scaling. The gate is computed from the
original equation rather than a router score or learned confidence and is re-evaluated after correction.

If $e_i$ fails, produces non-finite values, violates a contract, or is rejected by
$g_q$, the next candidate starts from the original request state. A dynamic step is
returned only after all state, derivative, event, and mode checks pass. Thus selection
error and candidate error affect runtime and solve completion, but not the acceptance rule.

\begin{proposition}[Verified-result return]
\label{prop:verified-return}
Assume that every result-return path invokes $g_q$ on the original request contract
and returns only on acceptance. Then any successfully returned state $y$ satisfies
$g_q(y)=\mathrm{accept}$, independently of the router ordering and of earlier
candidate, corrector, or device failures.
\end{proposition}

\begin{proof}
Each candidate path either fails, is rejected by $g_q$, or returns after acceptance.
The first two cases proceed to another stage and cannot produce a successful result.
The terminal solver is subject to the same rule. Therefore every successful return
has passed $g_q$.
\end{proof}

Proposition~\ref{prop:verified-return} is a numerical solver control-flow
invariant. It makes the acceptance boundary explicit but is not a theorem about an
incorrectly specified residual callback or faulty arithmetic hardware.

\subsection{Claims and Scope}

The formulation supports three testable claims. First, complete-cost selection can
outperform a fixed expert on qualified workloads. Second, correction can make an
otherwise unusable learned candidate profitable without allowing model output to
bypass original-equation acceptance. Third, mandatory original-equation gates reject observed bad
candidates and preserve the verified-result return invariant, while the experiments test
the implemented gates and continuation paths on observed workloads. The claim does not extend
to arbitrary incorrect callbacks or floating-point hardware faults. These claims concern
solver-internal numerical mechanisms and their observed experimental evidence.
